\section{Numerical experiments}
\label{sec:ct-experiments}

We compare minibatch and full-batch SFRAMES with
SPRING--SARAH~\cite{spring}, ProxSARAH~\cite{proxsarah}, and PALM~\cite{palm}
on low-rank reconstruction from real low-dose computed tomography (CT)
data. The experiment examines data fitting and constraint satisfaction at
equal projection budgets.

\paragraph{Reconstruction problem.}
The negative logarithm of calibrated X-ray transmission approximates line
integrals of attenuation. A \emph{view} groups detector rays at one source
angle. For an image \(X\), the linear forward map for view \(i\) is
\([P_i(X)]_j=\sum_{p,q}\ell_{ij,pq}X_{pq}\), where \(\ell_{ij,pq}\) is
the length of ray \(j\) through pixel \((p,q)\). Its adjoint \(P_i^\ast\)
backprojects detector residuals onto the image using the same weights.

We use low-dose fan-beam measurements from 2DeteCT~\cite{detect} and
represent the \(512\times512\) attenuation image as \(UV^\top\), with
\(U,V\in\mathbb R^{512\times64}\). For a set \(\mathcal I\) of \(N\)
training views and detector data \(y_i\in\mathbb R^d\), \(d=956\), the
training loss is
\begin{equation}
  f(U,V)=\frac1N\sum_{i\in\mathcal I}
          \frac{\|P_i(UV^\top)-y_i\|_2^2}{2d}.
  \label{eq:ct-loss}
\end{equation}
We constrain each factor by a spectral norm bound and require nonnegative
factors, which ensure nonnegative attenuation. For \(x=(U,V)\), set
\[
 C=\{(U,V):\|U\|_{\rm op}\le R_U,\ \|V\|_{\rm op}\le R_V\},
 \qquad D=\{(U,V):U\ge0,\ V\ge0\}.
\]
The bilinear factorization makes the problem nonconvex. Its SFRAMES
formulation is \(\min_{x\in C} f(x)+g(Tx)\), with
\(T=\operatorname{Id}\) and \(g=\iota_D\), the indicator of \(D\).
Uniformly sampling views gives an unbiased estimator of \(\nabla f\).
The spectral-ball linear oracle separates by factor and uses one thin SVD
per block. The gradient of the Moreau envelope of \(g\) is
\((x-x_+)/\beta\), where \(x_+\) clips both factors to nonnegativity.

\paragraph{Evaluation.}
We evaluate on 14 slices from six acquisition groups, with either \(360\)
or \(1{,}080\) training views and \(360\) held-out views per slice.
The two view counts change the finite-sum size while keeping the factor
dimensions fixed. We use three stochastic repetitions per slice and view
count, with a common initialization across methods.
We compare the last iterate within \(69{,}120\) forward and adjoint
views per direction, including all refresh and recursive-gradient work.
An oracle call returns one joint or block gradient on a batch.
All trajectories use observations recorded at every update.

Raw loss is the data-fit term \(f\) evaluated at the current factors.
To assess a feasible reconstruction, we clip each factor to
nonnegativity and rescale it radially into its spectral ball if needed.
We evaluate the resulting image by training and held-out loss and relative
Frobenius error against the high-dose reconstruction.
Repair is used only for evaluation. Raw nonnegativity violation is
\[
 d_D(U,V)=\bigl(\|U-U_+\|_F^2+\|V-V_+\|_F^2\bigr)^{1/2}.
\]
Relative infeasibility divides \(d_D\) by
\(\bigl(\|U_0\|_F^2+\|V_0\|_F^2\bigr)^{1/2}\).
Quality summaries take successive medians over repetitions, slices within
each group, and groups. Paired percentage changes are computed before this
aggregation, so they need not equal ratios of the table entries.

\paragraph{Results.}
@@RESULTS_TEXT@@
Table~\ref{tab:ct-results} gives endpoint metrics; Figures~\ref{fig:ct-trajectories-360}
and~\ref{fig:ct-trajectories-1080} show progress versus iterations, oracle calls,
and processed views.

@@RESULTS_TABLE@@

\paragraph{Implementation and parameter selection.}
All methods start from a nonnegative SVD-based factorization of a
sensitivity-normalized backprojection of training views, with its
amplitude fitted to those views. Set
\(R_U=2\|U_0\|_{\rm op}\) and \(R_V=2\|V_0\|_{\rm op}\).
Minibatch SFRAMES samples \(64\) views uniformly without replacement,
independently across iterations; full-batch SFRAMES uses all \(N\) views.
Both use \(\gamma_k=\gamma_0(k+1)^{-3/4}\),
\(\alpha_k=(k+1)^{-1/2}\) for gradient averaging, and
\(\beta_k=\beta_0(k+1)^{-1/4}\), with
\(\beta_0=c_\beta/\max(L_U,L_V)\) scaled by initial block-curvature estimates.
@@TUNING_TEXT@@

SPRING--SARAH alternates block updates with loopless SARAH estimates;
ProxSARAH uses joint updates with periodic snapshots; PALM alternates
full-gradient block updates. The stochastic baselines use \(64\)-view
gradient differences. Their proximal steps project onto each spectral
ball intersected with the nonnegative orthant, using practical tuned
step sizes. The \CTReproducibility{} specifies the data split,
preprocessing, search grids, and numerical solver settings.

\begin{figure}[!htbp]
 \centering
 \includegraphics[width=\linewidth]{trajectories_360.pdf}
 \caption{Trajectories for \(N=360\). Thin curves show individual
 runs at their recorded iterates; bold curves summarize work and metrics
 at the same iteration by successive medians over repetitions, slices,
 and groups. SFRAMES bands span the middle \(50\%\) of group medians.
 Training losses are relative to their initial values.
 Vertical markers indicate one quarter of the total view budget.}
 \label{fig:ct-trajectories-360}
\end{figure}

\begin{figure}[!htbp]
 \centering
 \includegraphics[width=\linewidth]{trajectories_1080.pdf}
 \caption{Trajectories for \(N=1{,}080\), summarized as in
 Figure~\ref{fig:ct-trajectories-360}. The three projection baselines
 coincide at zero orthant distance.}
 \label{fig:ct-trajectories-1080}
\end{figure}
